% id: 141-52
% book: 쎈B 대수 · 10 수학적 귀납법
% section: 유형 13 수학적 귀납법; 부등식의 증명
% figure: none
% answer: 
% answer_source: 
% variant_level: 0 (원본 전사)
% 이 파일은 build.mjs 가 items.json 에서 만든다. 고칠 때는 items.json 을 고친다.
\smash{\raisebox{4.5mm}{\dmkichul{쎈B 대수 141쪽 52번}}}
다음은 모든 자연수 $n$에 대하여 부등식\par\smallskip\dispeq{$\dfrac{1!+2!+3!+\cdots+n!}{(n+1)!}<\dfrac{2}{n+1}$\qquad $\cdots\cdots$ ㉠}\par\smallskip\noindent 가 성립함을 수학적 귀납법으로 증명한 것이다.\exprbox{\begin{minipage}{0.96\linewidth}\raggedright {\small\bfseries 증명}\\[1mm] (i) $n=1$일 때,\\[1mm] \hspace*{1.5em}(좌변)$=\dfrac{1!}{2!}=\dfrac{1}{2}$, (우변)$=\dfrac{2}{1+1}=1$\\[1mm] \hspace*{1em}따라서 ㉠이 성립한다.\\[1mm] (ii) $n=k$일 때 ㉠이 성립한다고 가정하면\\[1mm] \hspace*{1.5em}$\dfrac{1!+2!+3!+\cdots+k!}{(k+1)!}<\dfrac{2}{k+1}$\\[1mm] \hspace*{1em}$n=k+1$일 때\\[1mm] \hspace*{1.5em}$\dfrac{1!+2!+3!+\cdots+(k+1)!}{(k+2)!}$\\[1mm] \hspace*{2.5em}$=\fbox{(가)}\left\{\dfrac{1!+2!+3!+\cdots+k!}{(k+1)!}+\dfrac{(k+1)!}{(k+1)!}\right\}$\\[1mm] \hspace*{2.5em}$<\fbox{(가)}\left(1+\dfrac{2}{k+1}\right)$\\[1mm] \hspace*{2.5em}$=\dfrac{1}{k+2}+\fbox{(나)}$\\[1mm] \hspace*{1em}이때 자연수 $k$에 대하여 $\dfrac{2}{k+1}\le 1$이므로\\[1mm] \hspace*{1.5em}$\fbox{(나)}\le\dfrac{1}{k+2}$이고\\[1mm] \hspace*{1.5em}$\dfrac{1!+2!+3!+\cdots+(k+1)!}{(k+2)!}<\dfrac{2}{k+2}$이다.\\[1mm] \hspace*{1em}따라서 $n=k+1$일 때도 ㉠이 성립한다.\\[1mm] (i), (ii)에서 모든 자연수 $n$에 대하여 ㉠이 성립한다.\end{minipage}}\noindent 위의 과정에서 (가), (나)에 알맞은 식을 각각 $f(k)$, $g(k)$라 할 때, $f(8)+g(3)$의 값은?
\dmchoicesauto{\rule[-3ex]{0pt}{8ex}}{$\dfrac{1}{10}$}{$\dfrac{1}{5}$}{$\dfrac{3}{10}$}{$\dfrac{2}{5}$}{$\dfrac{1}{2}$}
